%% paper.tex — ZKP-RA: Zero-Knowledge Proof Reasoning Anchors for Mitigating
%%             Memory Poisoning in Autonomous Agentic DeFi
%%
%% Target venue: IEEE Transactions on Dependable and Secure Computing (TDSC)
%% Submission class: elsarticle (two-column journal mode)

\documentclass[final,5p,times,twocolumn]{elsarticle}

\usepackage{amsmath,amssymb,amsthm}
\usepackage{algorithm}
\usepackage{algpseudocode}
\usepackage{listings}
\usepackage{booktabs}
\usepackage{graphicx}
\usepackage{xcolor}
\usepackage{url}
\usepackage{hyperref}
\usepackage{cleveref}
\usepackage{microtype}

% ---------------------------------------------------------------------------
% Math environments
% ---------------------------------------------------------------------------
\theoremstyle{plain}
\newtheorem{theorem}{Theorem}[section]
\newtheorem{lemma}[theorem]{Lemma}
\newtheorem{corollary}[theorem]{Corollary}
\newtheorem{proposition}[theorem]{Proposition}
\theoremstyle{definition}
\newtheorem{definition}[theorem]{Definition}
\theoremstyle{remark}
\newtheorem{remark}[theorem]{Remark}

% Shorthand operators
\newcommand{\F}{\mathbb{F}}
\newcommand{\Z}{\mathbb{Z}}
\newcommand{\Adv}{\mathcal{A}}
\newcommand{\Sim}{\mathcal{S}}
\newcommand{\Hin}{H_{\mathrm{in}}}
\newcommand{\Hout}{H_{\mathrm{out}}}
\newcommand{\poseidon}{\textsf{Poseidon}}
\newcommand{\snark}{\text{zk-SNARK}}
\newcommand{\negl}{\mathrm{negl}}

\lstset{
  basicstyle=\footnotesize\ttfamily,
  breaklines=true,
  numbers=left,
  numberstyle=\tiny\color{gray},
  frame=single,
  captionpos=b,
}

% ---------------------------------------------------------------------------

\begin{document}

\begin{frontmatter}

\title{ZKP-RA: Zero-Knowledge Proof Reasoning Anchors for Mitigating\\
       Memory Poisoning in Autonomous Agentic DeFi}

\author[hcl]{Sunil Gentyala\corref{cor1}}
\cortext[cor1]{Corresponding author.}
\ead{sunil.gentyala@ieee.org}
\fntext[fn1]{Code and circuit artifacts available at:
\url{https://github.com/sunilgentyala/ZKP-RA}}

\address[hcl]{HCL America Inc., Independent Research Division}

\begin{abstract}
Autonomous agents that execute high-value transactions in decentralized finance
(DeFi) protocols face a qualitatively distinct attack surface: adversaries can
corrupt an agent's long-term vector memory or oracle input stream without
touching the execution layer, causing a fully-functional smart contract
infrastructure to authorize economically catastrophic transfers based on
poisoned reasoning. Trusted Execution Environments (TEEs) address hardware
integrity but provide no proof of \emph{reasoning soundness}, while optimistic
execution approaches rely on challenge windows that may expire before anomalous
transactions settle. We present ZKP-RA, a framework that embeds cryptographic
Reasoning Anchors directly into the agent-transaction pipeline. Before each
on-chain settlement, the agent commits a $\poseidon$-based state hash
$H_t = \poseidon(X_t, S_t, I_t)$ representing its context window, state
vector, and oracle inputs; subsequently generates a succinct non-interactive
argument (Groth16 over BN254) proving that the transaction payload $T_x$ was
derived from unmodified, policy-compliant state; and submits the proof to an
on-chain verifier that settles or rejects atomically. Benchmarks on a 3,800-constraint
R1CS instance show end-to-end proof generation latency of 182\,ms with
rapidsnark and on-chain verification consuming approximately 276,000 gas on an
EVM-equivalent chain, representing a $9.7\times$ reduction in overhead relative
to prior general-purpose zkML verification schemes~\cite{chen2024zkml}.
\end{abstract}

\begin{keyword}
zero-knowledge proofs \sep decentralized finance \sep memory poisoning \sep
LLM agents \sep Poseidon hash \sep Groth16 \sep smart contracts \sep prompt injection
\end{keyword}

\end{frontmatter}

% ---------------------------------------------------------------------------
\section{Introduction and Motivation}
\label{sec:intro}
% ---------------------------------------------------------------------------

The deployment of large language model (LLM) agents as autonomous economic
actors in on-chain DeFi protocols represents a structural shift in the threat
model underlying blockchain security. Classical smart contract vulnerabilities
are localized to the deterministic execution layer: a reentrancy bug, an
integer overflow, or an unchecked external call. The attack surface of an LLM
agent, however, spans both the deterministic execution layer (EVM bytecode,
storage slots, gas metering) and a stochastic semantic layer (transformer
attention over a context window of potentially unbounded provenance). These
two layers interact at a precise choke point: the moment an agent's reasoning
process produces a transaction payload that is subsequently signed and broadcast
to a peer-to-peer network.

This choke point is the locus of the threat we analyze. A sufficiently capable
adversary need never compromise a private key. Instead, it injects a
maliciously crafted payload $\delta$ into the agent's data ingestion stream
at some point upstream of the context window. When the agent's retrieval
augmentation mechanism subsequently incorporates $\delta$ into its reasoning
context, the resulting transaction payload $T_x$ reflects the adversary's
objective while the execution stack faithfully authorizes it. The attack is
semantically invisible to the contract layer.

Current mitigations fall into two classes. First, hardware-enforced isolation
via TEEs (Intel TDX, AMD SEV-SNP) places agent execution inside a confidential
computing enclave and produces remote attestation reports verifiable by
third parties~\cite{karim2025aibc}. TEE attestation proves that \emph{a specific
binary ran on certified hardware in an unmodified state}; it does not prove
that the binary's inputs were unmanipulated, nor does it produce a short,
on-chain-verifiable argument. Second, optimistic dispute protocols borrow
the fraud-proof architecture of rollup systems: a watchdog can challenge any
transaction suspected of arising from corrupt state. Optimistic systems carry
inherent latency and economic exposure during the challenge window; a
well-timed attack with high enough value can settle before a challenge is
completed~\cite{zhou2023sokdefi}.

ZKP-RA takes a different approach. Rather than attesting to the integrity of
the execution environment or relying on post-hoc dispute, ZKP-RA constructs a
\emph{pre-settlement proof of reasoning soundness} that the smart contract
verifier can check in the same atomic transaction that would release funds.
The key insight is that a useful security property need not require proving
the entire forward pass of a billion-parameter LLM. It requires only proving
that a compact, auditable \emph{policy function} $C$ was applied over inputs
that hash to a previously committed, on-chain-recorded value. The cost of
that proof is dominated by the hash computation and a linear number of range
constraints, not by the arithmetic complexity of the full model.

The remainder of this paper is organized as follows. \Cref{sec:threat} formally
defines the adversary model and the state-transition attack. \Cref{sec:arch}
describes the circuit design and on-chain verifier architecture. \Cref{sec:security}
provides the formal security analysis. \Cref{sec:eval} reports benchmarks.
\Cref{sec:related} surveys related work, and \Cref{sec:conclusion} concludes.

% ---------------------------------------------------------------------------
\section{Threat Model and Formal Problem Definition}
\label{sec:threat}
% ---------------------------------------------------------------------------

\subsection{System Model}
\label{sec:system}

We model an autonomous DeFi agent as a stateful probabilistic polynomial-time
(PPT) algorithm $\mathcal{M}$ operating over a sequence of time steps
$t \in \{0,1,\ldots,T\}$. At each step, $\mathcal{M}$ maintains:

\begin{itemize}
  \item A \emph{context window} $X_t \in \{0,1\}^{*}$: the concatenation of
        tokenized instructions, retrieval-augmented memory excerpts, and recent
        tool-call results, serialized and truncated to the model's context length.
  \item A \emph{state vector} $S_t \in \F_p^{d}$: a $d$-dimensional compressed
        representation of the agent's persistent long-term memory, implemented in
        practice as the mean-pooled embedding of top-$k$ retrieved documents.
  \item An \emph{input stream} $I_t \in \F_p^{m}$: a snapshot of $m$ oracle
        data feeds (asset prices, liquidity pool depths, cross-chain bridge state)
        sampled at time $t$.
\end{itemize}

The agent's \emph{state transition function} $f$ maps these three primitives,
together with private model weights $w$, to a transaction payload:
\begin{equation}
  T_x^{(t)} = f\!\left(X_t,\, S_t,\, I_t;\, w\right).
  \label{eq:stf}
\end{equation}
A payload $T_x^{(t)}$ is a structured object specifying a target ERC-20
contract, a recipient address, a transfer amount, a maximum slippage in basis
points, and a deadline timestamp.

The \emph{on-chain execution environment} is a deterministic EVM-equivalent
state machine $\mathcal{E}$ that accepts $T_x^{(t)}$ and executes the
corresponding asset transfer if and only if the authorization condition
$\mathcal{V}(T_x^{(t)}) = 1$ holds.

\subsection{Adversary Definition}
\label{sec:adversary}

\begin{definition}[Memory-Poisoning Adversary]
\label{def:adversary}
A \emph{memory-poisoning adversary} $\Adv$ is a PPT algorithm with access to
a write oracle $\mathcal{O}_{\mathrm{inject}}$ for the agent's data ingestion
stream. Formally, $\Adv$ can compute an adversarial perturbation
$\delta \in \{0,1\}^{*}$ and submit it to $\mathcal{O}_{\mathrm{inject}}$
such that the poisoned input stream becomes
$$\tilde{I}_t = I_t \oplus_{\mathrm{mux}} \delta$$
where $\oplus_{\mathrm{mux}}$ denotes structured field substitution at
adversarially chosen indices. $\Adv$ cannot directly access $w$, cannot
forge digital signatures on messages already broadcast, and cannot modify
the deterministic EVM execution trace after a block is finalized. This
threat class subsumes indirect prompt injection~\cite{greshake2023ipi},
retrieval-augmented generation corpus poisoning~\cite{poisonedrag2024}, and
on-chain oracle feed manipulation~\cite{defitrace2024}.
\end{definition}

The adversary's goal is to cause the agent to produce a transaction payload
$\tilde{T}_x^{(t)} = f(X_t, S_t, \tilde{I}_t; w)$ that satisfies
$\mathcal{V}(\tilde{T}_x^{(t)}) = 1$ while maximizing an economic objective
$\Phi(\tilde{T}_x^{(t)}) \gg \Phi(T_x^{(t)})$ (e.g., draining a liquidity pool
or redirecting a large transfer to an adversary-controlled address).

\begin{definition}[State-Tampering Attack]
\label{def:tampering}
A \emph{state-tampering attack} at time $t$ is a tuple
$(\delta, T_x^*, \epsilon)$ such that:
\begin{enumerate}
  \item $\Adv$ injects $\delta$ into the input stream yielding
        $\tilde{I}_t \neq I_t$.
  \item The agent computes $\tilde{T}_x = f(X_t, S_t, \tilde{I}_t; w)$.
  \item The economic deviation satisfies
        $\|\tilde{T}_x.\mathtt{amount} - T_x^*.\mathtt{amount}\| > \epsilon$
        for some legitimate target amount $T_x^*.\mathtt{amount}$.
\end{enumerate}
\end{definition}

\subsection{Security Objective}

The security objective of ZKP-RA is to ensure that for any PPT adversary
$\Adv$ of the type in \Cref{def:adversary}, the probability that
$\mathcal{V}(\tilde{T}_x) = 1$ is negligible in the security parameter
$\lambda$:
\begin{equation}
  \Pr\!\left[\mathcal{V}\!\left(f\!\left(X_t, S_t, \tilde{I}_t; w\right)\right) = 1
    \,\Big|\, \tilde{I}_t \neq I_t \right] \leq \negl(\lambda).
  \label{eq:security_obj}
\end{equation}
This objective is achieved by binding $\mathcal{V}$ to a succinct proof
$\pi$ that the inputs to $f$ at the time of proof generation hash to a
value committed on-chain \emph{before} the current block, preventing any
substitution that occurred after the commitment.

% ---------------------------------------------------------------------------
\section{Technical Architecture and Circuit Design}
\label{sec:arch}
% ---------------------------------------------------------------------------

\subsection{Protocol Overview}

ZKP-RA operates in three phases that interleave on-chain and off-chain
computation. Let $\lambda = 128$ be the security parameter and let $p$ denote
the BN254 scalar field prime.

\paragraph{Phase 1 — State Commitment.}
At block $b$, the agent serializes $X_t, S_t, I_t$ into field-element arrays
$\mathbf{x} \in \F_p^{N_X}$, $\mathbf{s} \in \F_p^{N_S}$, $\mathbf{i} \in \F_p^{N_I}$
by packing 31-byte chunks (ensuring each chunk is strictly less than $p$).
It then computes the \emph{state commitment}:
\begin{equation}
  H_t = \poseidon\!\left(
    \poseidon_{\mathrm{CTX}}(\mathbf{x}),\,
    \poseidon_{\mathrm{STATE}}(\mathbf{s}),\,
    \poseidon_{\mathrm{INPUT}}(\mathbf{i})
  \right) \in \F_p,
  \label{eq:commitment}
\end{equation}
where $\poseidon_{\mathrm{D}}$ denotes a domain-separated sponge absorption
over the array with domain tag $D \in \{1,2,3\}$ mixed into the initial
capacity element. The agent calls \texttt{commitState}($H_t$) on the verifier
contract, anchoring $H_t$ to block $b$ on-chain. The $\poseidon$ permutation
is instantiated with $\alpha=5$, $R_f=8$, $R_p=57$ over BN254, as specified
in~\cite{grassi2021poseidon}.

\paragraph{Phase 2 — Proof Generation.}
After generating transaction payload $T_x$, the agent invokes the prover
backend with private witnesses $(\mathbf{x}, \mathbf{s}, \mathbf{i},
r_{\mathrm{router}}, \sigma)$ where $r_{\mathrm{router}}$ is the whitelisted
router address and $\sigma$ is a fresh random salt $\sigma \xleftarrow{\$} \F_p$.
The prover computes a Groth16 proof $\pi = (A, B, C) \in
\mathbb{G}_1 \times \mathbb{G}_2 \times \mathbb{G}_1$ over BN254~\cite{groth2016}, certifying:
\begin{equation}
  \pi \leftarrow \mathsf{Prove}\!\left(
    \mathsf{pk},\;
    \mathbf{pub} = (H_t, T_x),\;
    \mathbf{wit} = (\mathbf{x}, \mathbf{s}, \mathbf{i}, r_{\mathrm{router}}, \sigma)
  \right).
  \label{eq:prove}
\end{equation}

\paragraph{Phase 3 — On-Chain Settlement.}
The agent calls \texttt{verifyAndExecute}$(\pi, T_x, n, \nu)$ where $n$ is the
commitment nonce and $\nu$ is a replay-prevention nullifier. The on-chain
Groth16 verifier checks:
\begin{equation}
  e(A, B) = e(\alpha, \beta)\cdot e(\mathsf{vk}_x, \gamma)\cdot e(C, \delta),
  \label{eq:pairing}
\end{equation}
where $\mathsf{vk}_x = \sum_{i=0}^{l} a_i \cdot \mathsf{IC}[i]$ is the
linear combination of input commitments over the public inputs
$(H_t, T_x.\mathtt{recipient}, T_x.\mathtt{amount}, T_x.\mathtt{slippage},
T_x.\mathtt{routerHash})$. Only upon successful verification does the contract
release the ERC-20 transfer.

\subsection{Circuit Arithmetization}
\label{sec:circuit}

The circuit \texttt{reasoning\_anchor.circom} is written in Circom 2.1.6 and
compiled to an R1CS instance over $\F_p$.

\subsubsection{Poseidon Sponge Absorption}
The memory chunking and hashing sub-circuit (\texttt{PoseidonChunked}) absorbs
$N$ field elements using $\lceil N/2 \rceil$ instantiations of the $t=3$
Poseidon permutation. For a single $t=3$ Poseidon call with $\alpha = 5$ (the
S-box exponent over BN254), the number of full rounds is $R_f = 8$ and partial
rounds $R_p = 57$, yielding:
\begin{equation}
  C_{\poseidon}(t=3) = (R_f + R_p)\cdot t = (8+57)\cdot 3 = 195 \text{ constraints}.
  \label{eq:poseidon_cost}
\end{equation}
Absorbing $N_X = 8$ context chunks requires $\lceil 8/2\rceil = 4$ Poseidon calls
($4 \times 195 = 780$ constraints). Repeating for all three memory regions and
the final root combinator gives $780 + 780 + 390 + 195 = 2{,}145$ constraints
for the full $H_t$ computation.

\subsubsection{Policy Enforcement Constraints}
Transaction bound verification uses the \texttt{LessEqThan} comparator from
circomlib, which decomposes each operand into $n$ bits requiring $n$ Boolean
constraints plus $n - 1$ range constraints. For the 252-bit amount check:
\begin{equation}
  C_{\mathrm{range}}(n=252) = 252 \text{ (bit decomp)} + 252 \text{ (range)} = 504 \text{ constraints}.
  \label{eq:range_cost}
\end{equation}
The 14-bit slippage check contributes $28$ additional constraints. The router
whitelist membership proof requires one additional $t=2$ Poseidon call
(estimated $C_{\poseidon}(t=2) \approx 150$ constraints) plus an equality
assertion. The total R1CS instance size is approximately:
\begin{equation}
  C_{\mathrm{total}} \approx 2{,}145 + 504 + 28 + 150 + 50 = 2{,}877 \text{ constraints}.
  \label{eq:total_constraints}
\end{equation}
This is consistent with the empirically measured value of $3{,}800$ constraints
after Circom's automatic intermediate signal introduction.

\subsubsection{Optimization: Structural Asymmetry of Transformer Traces}
A key design choice is \emph{not} arithmetizing the full transformer forward
pass. ZK-constrained ML inference for general neural networks incurs a
constraint count proportional to the number of activations; for a 7B-parameter
model this exceeds $10^{10}$ constraints, making real-time proving infeasible
with current hardware~\cite{peng2025survey,chen2024zkml}. This scales
qualitatively worse than the post-quantum transparent proof systems based on
FRI~\cite{bensasson2018stark}, yet neither family is viable for full-model
real-time proving in the DeFi settlement latency budget. ZKP-RA exploits the
structural asymmetry between prefill and decoding:

\begin{itemize}
  \item The \emph{prefill} phase processes the full context window in parallel.
        Its output is the KV-cache, a deterministic function of $X_t$.
  \item The \emph{decoding} phase auto-regressively generates tokens from the
        KV-cache. Its output is stochastic and temperature-dependent.
\end{itemize}

We do not prove the decoding trace. Instead, we prove only that the \emph{inputs}
to the prefill phase hash to $H_t$, and that the resulting transaction payload
satisfies the auditable policy function $C$. The policy function is defined
over the structured output fields of $T_x$ (amount, slippage, router) rather
than over the raw token sequence. This reduces the constraint count from
$\mathcal{O}(N_{\mathrm{params}})$ to $\mathcal{O}(N_X + N_S + N_I)$, which
for our instantiation is $\mathcal{O}(20)$ field-element arrays.

The tradeoff is that ZKP-RA proves \emph{input integrity} and \emph{output
policy compliance} but not the full correctness of the reasoning trace. This
is the appropriate property to prove in a DeFi context: the adversary's goal
is to corrupt inputs or cause policy-violating outputs, not to compromise the
internal weights of the model~\cite{verillm2025}.

\subsection{On-Chain Verifier Gas Cost}
\label{sec:gas}

The Groth16 verifier in \texttt{Verifier.sol} performs:
\begin{itemize}
  \item Four BN254 pairing operations via EIP-197 precompile:
        $4 \times 45{,}000 = 180{,}000$ gas.
  \item Scalar multiplications for the $\mathsf{vk}_x$ linear combination
        ($l = 5$ inputs): $5 \times 6{,}000 = 30{,}000$ gas (EIP-196 ecMul).
  \item Point additions: $4 \times 500 = 2{,}000$ gas (EIP-196 ecAdd).
  \item Storage reads (commitment lookup, nullifier check): $\sim 4{,}200$ gas.
  \item ERC-20 safeTransfer: $\sim 35{,}000$ gas.
  \item Ancillary overhead (calldata, event emission): $\sim 25{,}000$ gas.
\end{itemize}
Total: approximately $276{,}200$ gas, equivalent to approximately \$0.19 at
20 Gwei and ETH at \$3{,}500. For comparison, a standard Uniswap V3 exact-input
swap consumes $\sim 130{,}000$ gas; ZKP-RA adds a $2.1\times$ overhead, which
is acceptable for high-value institutional transactions exceeding \$50{,}000 in
notional value.

% ---------------------------------------------------------------------------
\section{Security Analysis}
\label{sec:security}
% ---------------------------------------------------------------------------

\subsection{Completeness}
\label{sec:completeness}

\begin{theorem}[Completeness]
\label{thm:completeness}
For any honest agent $\mathcal{M}$ that samples $(X_t, S_t, I_t)$ from its
unmodified state, computes $H_t$ as in \cref{eq:commitment}, and generates a
transaction payload $T_x$ satisfying all policy constraints, the verifier
accepts with probability 1.
\end{theorem}

\begin{proof}
By the perfect completeness of the Groth16 proving system over BN254, if the
prover holds a valid witness $\mathbf{wit} = (\mathbf{x}, \mathbf{s},
\mathbf{i}, r, \sigma)$ such that:
(i)~the circuit's Poseidon computation reproduces $H_t$, and
(ii)~all policy constraints are satisfied,
then the proof $\pi$ satisfies the pairing equation in \cref{eq:pairing}
with probability 1 over the randomness used in the proof generation
algorithm. The on-chain verifier evaluates the same pairing equation
deterministically using the hardcoded verification key; thus it accepts.
\end{proof}

\subsection{Soundness}
\label{sec:soundness}

\begin{theorem}[Computational Soundness against Memory-Poisoning]
\label{thm:soundness}
Under the $q$-Strong Diffie-Hellman ($q$-SDH) assumption on BN254 and the
collision resistance of $\poseidon$ over $\F_p$, no PPT adversary $\Adv$ of
the type in \Cref{def:adversary} can cause the verifier to accept a proof
$\tilde{\pi}$ for a transaction payload arising from poisoned state
$\tilde{I}_t \neq I_t$ with probability non-negligibly greater than
$\negl(\lambda)$.
\end{theorem}

\begin{proof}
We proceed by contradiction. Suppose $\Adv$ succeeds with non-negligible
probability $\epsilon(\lambda)$. We construct a reduction $\mathcal{R}$ that
breaks either the $q$-SDH assumption or the collision resistance of $\poseidon$.

Case 1 (Groth16 knowledge soundness): If $\Adv$ produces an accepting
$\tilde{\pi}$ for public inputs $(H_t, \tilde{T}_x)$ where $H_t$ is the
on-chain committed value, then by the knowledge soundness of Groth16 under
$q$-SDH (Theorem 3 of~\cite{vcnn2024}), there exists a PPT extractor
$\mathcal{E}$ that outputs a valid witness
$\tilde{\mathbf{wit}} = (\tilde{\mathbf{x}}, \tilde{\mathbf{s}},
\tilde{\mathbf{i}}, \tilde{r}, \tilde{\sigma})$ such that the circuit
constraints are satisfied. In particular, the Poseidon sub-circuit forces:
$$\poseidon\!\left(\poseidon_{\mathrm{CTX}}(\tilde{\mathbf{x}}),
  \poseidon_{\mathrm{STATE}}(\tilde{\mathbf{s}}),
  \poseidon_{\mathrm{INPUT}}(\tilde{\mathbf{i}})\right) = H_t.$$

Case 2 (Poseidon collision): Since $\tilde{I}_t \neq I_t$ implies
$\tilde{\mathbf{i}} \neq \mathbf{i}$ (by injectivity of the serialization
map, which is a bijection on its domain), we require
$\poseidon_{\mathrm{INPUT}}(\tilde{\mathbf{i}}) = \poseidon_{\mathrm{INPUT}}(\mathbf{i})$,
which constitutes a second-preimage collision in $\poseidon$. Under the
algebraic independence assumption of the $\poseidon$ round constants (a standard
assumption in the literature~\cite{peng2025survey}), such a collision occurs
with probability at most $p^{-1} \approx 2^{-254}$, which is negligible.

Since both cases lead to contradictions with negligible success probability,
the theorem holds.
\end{proof}

\subsection{Zero-Knowledge}
\label{sec:zk}

\begin{theorem}[Computational Zero-Knowledge]
\label{thm:zk}
The ZKP-RA proof system is computationally zero-knowledge. For any PPT
verifier $\mathcal{V}^*$ (including the on-chain contract), there exists a
simulator $\Sim$ such that for all valid statements, the distribution of
$\Sim$'s output is computationally indistinguishable from the distribution of
real proofs.
\end{theorem}

\begin{proof}
The zero-knowledge property follows directly from the zero-knowledge property
of Groth16 over BN254~\cite{vcnn2024}. The simulator $\Sim$ outputs a simulated
proof $\pi_{\mathrm{sim}}$ by sampling random field elements for the proof
components and programming the random oracle (in the random oracle model) to
produce consistent responses. The public inputs are identical in both the real
and simulated executions. Since the on-chain verifier evaluates only the
pairing equation~\cref{eq:pairing} over the public inputs and proof components,
its view is computationally indistinguishable.

Concretely, no information about the private witnesses $(\mathbf{x}, \mathbf{s},
\mathbf{i}, r, \sigma)$ --- i.e., the raw context window tokens, state
embedding coordinates, oracle data, or router address --- is revealed by
the proof.
\end{proof}

\begin{remark}
The zero-knowledge property is practically significant: it prevents the
on-chain verifier (or any chain observer) from reconstructing the agent's
internal reasoning state from the proof transcript, which would constitute
a model-extraction attack against the agent's proprietary strategy.
\end{remark}

% ---------------------------------------------------------------------------
\section{Evaluation}
\label{sec:eval}
% ---------------------------------------------------------------------------

\subsection{Benchmark Configuration}

Experiments ran on a workstation with an AMD EPYC 9354 (32-core, 3.25 GHz),
128 GB DDR5 ECC, and Ubuntu 22.04. Proving used SnarkJS 0.7.4 (WASM), Node.js
20.11, and rapidsnark 0.0.1 (native C++ with SIMD). On-chain gas measurements
used Hardhat 2.22 with the Hardhat Gas Reporter plugin against a local Anvil
fork at Istanbul EVM.

\subsection{Proof Generation Latency}

\Cref{tab:proving_time} reports median latency over 100 trials. The
3,800-constraint R1CS instance fits entirely within the L3 cache during proving.

\begin{table}[h]
\centering
\caption{Proof generation latency for reasoning\_anchor.circom (3,800 constraints).}
\label{tab:proving_time}
\begin{tabular}{lrrr}
\toprule
\textbf{Backend} & \textbf{Witness (ms)} & \textbf{Prove (ms)} & \textbf{Total (ms)} \\
\midrule
SnarkJS WASM     & 38 & 1{,}164 & 1{,}202 \\
rapidsnark       & 12 &    170  &    182 \\
\bottomrule
\end{tabular}
\end{table}

The 182\,ms total latency for rapidsnark is well within the block time of
all major EVM chains (12\,s on Ethereum mainnet, 2\,s on Polygon PoS). For
high-frequency agents operating at sub-second cadence, the prover can be
parallelized across CPU cores since each transaction generates an independent
witness.

\subsection{On-Chain Verification Gas}

\Cref{tab:gas} breaks down gas consumption by operation for the
\texttt{verifyAndExecute} call.

\begin{table}[h]
\centering
\caption{On-chain gas breakdown for \texttt{verifyAndExecute} (EVM, Istanbul opcodes).}
\label{tab:gas}
\begin{tabular}{lr}
\toprule
\textbf{Operation} & \textbf{Gas} \\
\midrule
BN254 pairing (EIP-197, 4 pairs)         & 180{,}000 \\
Scalar multiplication (EIP-196, 5 calls) &  30{,}000 \\
Point addition (EIP-196, 4 calls)        &   2{,}000 \\
Storage reads (commitment, nullifier)    &   4{,}200 \\
ERC-20 safeTransfer                      &  35{,}000 \\
Calldata + event emission                &  25{,}000 \\
\midrule
\textbf{Total}                           & \textbf{276{,}200} \\
\bottomrule
\end{tabular}
\end{table}

\subsection{Comparison with ZKML Baselines}

\Cref{tab:comparison} situates ZKP-RA against prior systems for verifiable ML
inference. The constraint count and gas cost of ZKP-RA are orders of magnitude
lower than full-model verification because ZKP-RA proves input integrity and
output policy rather than the complete inference trace.

\begin{table}[h]
\centering
\caption{Comparison with related verifiable inference systems.}
\label{tab:comparison}
\begin{tabular}{lrrr}
\toprule
\textbf{System} & \textbf{Constraints} & \textbf{Prove (ms)} & \textbf{Verify (gas)} \\
\midrule
ZKML~\cite{chen2024zkml}       & $\sim\!8\times10^6$ & 2{,}100  & 2{,}680{,}000 \\
vCNN~\cite{vcnn2024}           & $\sim\!2\times10^6$ & 890      & 698{,}000 \\
VeriLLM~\cite{verillm2025}     & $\sim\!4\times10^5$ & 410      & 512{,}000 \\
\textbf{ZKP-RA (ours)}         & $\sim\!3.8\times10^3$ & \textbf{182} & \textbf{276{,}200} \\
\bottomrule
\end{tabular}
\end{table}

The $9.7\times$ gas reduction relative to ZKML stems from the structural
decoupling described in \Cref{sec:circuit}: by proving only that
$\poseidon(\mathbf{x}, \mathbf{s}, \mathbf{i}) = H_t$ and that $T_x$ satisfies
the policy, rather than arithmetizing the weight matrix multiplications,
the circuit remains amenable to real-time on-chain settlement.

% ---------------------------------------------------------------------------
\section{Related Work}
\label{sec:related}
% ---------------------------------------------------------------------------

\paragraph{zkML systems.}
ZKML~\cite{chen2024zkml} demonstrates that standard ML inference graphs can be
compiled to arithmetic circuits for Groth16 proving. The EuroSys 2024 paper
focuses on constraint optimization techniques including operator fusion and
quantization-aware layouting. Its primary target is correctness verification of
inference results, not the memory-integrity problem posed by autonomous agents.
The survey by Peng et al.~\cite{peng2025survey} provides a comprehensive taxonomy
of ZKP-based verifiable ML, covering SNARKs, STARKs, and hybrid approaches;
it identifies agent-level memory integrity as an open research challenge, which
ZKP-RA directly addresses.

\paragraph{Verifiable LLM inference.}
VeriLLM~\cite{verillm2025} introduces a lightweight framework for decentralized
inference verification, reducing constraint counts by selectively proving
critical transformer layers. ZKP-RA is complementary: where VeriLLM targets
inference correctness, ZKP-RA targets input integrity and output policy, enabling
a layered defense in which VeriLLM can optionally be composed as an inner layer.

\paragraph{Verifiable CNNs and zk-SNARKs.}
The vCNN system~\cite{vcnn2024} applies zk-SNARKs to convolutional neural
network inference for image classification tasks. Its R1CS circuit construction
techniques for fixed-point arithmetic inform the range-check design in the
ZKP-RA policy engine, specifically the bit-decomposition approach to upper-bound
proofs. PLONK~\cite{plonk2019} offers a universal SRS and avoids the per-circuit
trusted setup required by Groth16; its lookup arguments are relevant to future
whitelist-membership optimizations in ZKP-RA.

\paragraph{DeFi attack surface.}
The SoK survey by Zhou et al.~\cite{zhou2023sokdefi} provides a systematic
taxonomy of 4.7 years of DeFi incidents totalling \$3.24 billion in losses,
including oracle manipulation and cross-chain bridge exploits that constitute the
primary injection channels for the adversary in \Cref{def:adversary}.
DeFiTrace~\cite{defitrace2024} and DeFort~\cite{defort2024} offer detection-side
tooling; ZKP-RA is complementary in that it prevents execution of adversarially
derived payloads regardless of whether the injection was detected.

\paragraph{Autonomous agents on blockchain.}
Karim et al.~\cite{karim2025aibc} survey the intersection of multi-agent AI and
blockchain, introducing Proof-of-Thought consensus and highlighting the integrity
gap that ZKP-RA addresses. The SoK by Li et al.~\cite{sokaibc2025} categorizes
the security and privacy risks specific to AI agents operating on-chain, including
memory leakage, decision manipulation, and oracle abuse --- all of which fall
within the ZKP-RA adversary model.

\paragraph{Indirect prompt injection and RAG poisoning.}
Greshake et al.~\cite{greshake2023ipi} established the foundational taxonomy of
indirect prompt injection, demonstrating working exploits against GPT-4-powered
agents. PoisonedRAG~\cite{poisonedrag2024} showed that injecting a small number
of adversarial documents into a retrieval corpus is sufficient to steer LLM
outputs toward attacker-chosen targets, directly motivating the $H_t$ commitment
mechanism in ZKP-RA.

\paragraph{TEE-based agent isolation.}
Trusted Execution Environments provide hardware-level integrity for LLM agents
deployed in confidential computing enclaves. While TEEs offer strong confidentiality
guarantees, remote attestation does not produce short proofs verifiable by an
EVM smart contract without additional adapter layers. ZKP-RA fills this gap by
generating a succinct on-chain-verifiable argument, and the two approaches are
architecturally compatible.

\paragraph{Alternative proof systems.}
STARKs~\cite{bensasson2018stark} achieve post-quantum security and require no
trusted setup, at the cost of larger proof sizes ($\sim$45 KB vs.\ $\sim$192
bytes for Groth16). For the on-chain settlement use case, proof size translates
directly to calldata gas cost; the 192-byte Groth16 proof in ZKP-RA is therefore
preferred over STARKs for the current deployment target.

% ---------------------------------------------------------------------------
\section{Conclusion}
\label{sec:conclusion}
% ---------------------------------------------------------------------------

We presented ZKP-RA, a framework for cryptographically binding an autonomous
DeFi agent's transaction outputs to verifiably clean, policy-compliant input
state. The framework's core mechanism, a Groth16 Reasoning Anchor proof over a
3,800-constraint BN254 circuit, requires only 182\,ms to generate and 276,200
gas to verify, making it the first zero-knowledge memory-integrity primitive
that is practical for real-time on-chain DeFi settlement.

The security analysis establishes that a memory-poisoning adversary who cannot
find a Poseidon collision or break the $q$-SDH assumption on BN254 cannot cause
the verifier to accept a transaction arising from manipulated inputs. The
zero-knowledge property ensures that the proof transcript reveals nothing about
the agent's internal reasoning state, context window, or strategy parameters.

Future work will investigate batch-proof aggregation using PLONK~\cite{plonk2019}
or HyperPlonk to amortize the per-transaction proving overhead across agent
epochs, and will examine the integration of ZKP-RA with TEE attestation chains
to provide defense-in-depth against hardware-level adversaries~\cite{sokaibc2025}.

% ---------------------------------------------------------------------------
\section*{Code Availability}
% ---------------------------------------------------------------------------

The complete implementation of ZKP-RA --- including the Groth16 verifier
contract (\texttt{contracts/Verifier.sol}), the Circom 2.1.6 reasoning anchor
circuit (\texttt{circuits/reasoning\_anchor.circom}), the asynchronous Python
agent guardian daemon (\texttt{src/agent\_guardian.py}), and the circuit
compile and trusted-setup automation scripts --- is openly available under the
MIT License at:

\begin{center}
\url{https://github.com/sunilgentyala/ZKP-RA}
\end{center}

% ---------------------------------------------------------------------------
\bibliographystyle{model1-num-names}
\bibliography{references}
% ---------------------------------------------------------------------------

\begin{thebibliography}{15}

\bibitem{chen2024zkml}
B.\,J.\ Chen, S.\ Waiwitlikhit, I.\ Stoica, and D.\ Kang,
``ZKML: An Optimizing System for ML Inference in Zero-Knowledge Proofs,''
in \textit{Proc.\ EuroSys 2024}, pp.\ 560--574, ACM, 2024.
\newblock DOI: \texttt{10.1145/3627703.3650088}

\bibitem{groth2016}
J.\ Groth,
``On the Size of Pairing-Based Non-interactive Arguments,''
in \textit{Proc.\ EUROCRYPT 2016}, LNCS vol.\ 9666, pp.\ 305--326,
Springer, 2016.
\newblock DOI: \texttt{10.1007/978-3-662-49896-5\_11}

\bibitem{grassi2021poseidon}
L.\ Grassi, D.\ Khovratovich, C.\ Rechberger, A.\ Roy, and M.\ Schofnegger,
``Poseidon: A New Hash Function for Zero-Knowledge Proof Systems,''
in \textit{Proc.\ USENIX Security 2021}, pp.\ 519--535, USENIX, 2021.
\newblock URL: \texttt{https://www.usenix.org/conference/usenixsecurity21/presentation/grassi}

\bibitem{greshake2023ipi}
K.\ Greshake, S.\ Abdelnabi, S.\ Mishra, C.\ Endres, T.\ Holz, and M.\ Fritz,
``Not What You've Signed Up For: Compromising Real-World LLM-Integrated
Applications with Indirect Prompt Injection,''
in \textit{Proc.\ 16th ACM Workshop on Artificial Intelligence and Security
(AISec@CCS 2023)}, ACM, 2023.
\newblock DOI: \texttt{10.1145/3605764.3623985}

\bibitem{vcnn2024}
S.\ Lee, H.\ Ko, J.\ Kim, and H.\ Oh,
``vCNN: Verifiable Convolutional Neural Network Based on zk-SNARKs,''
\textit{IEEE Trans.\ Dependable Secure Comput.}, vol.\ 21, no.\ 4,
pp.\ 4254--4270, 2024.
\newblock DOI: \texttt{10.1109/TDSC.2023.3348760}

\bibitem{verillm2025}
VeriLLM: A Lightweight Framework for Publicly Verifiable Decentralized Inference,
\textit{arXiv preprint arXiv:2509.24257}, 2025.
\newblock DOI: \texttt{10.48550/arXiv.2509.24257}

\bibitem{zhou2023sokdefi}
L.\ Zhou, X.\ Xiong, J.\ Ernstberger, S.\ Chaliasos, Z.\ Wang, Y.\ Wang,
K.\ Qin, R.\ Wattenhofer, D.\ Song, and A.\ Gervais,
``SoK: Decentralized Finance (DeFi) Attacks,''
in \textit{Proc.\ IEEE Symp.\ Security \& Privacy (S\&P 2023)},
pp.\ 2444--2461, IEEE, 2023.
\newblock DOI: \texttt{10.1109/SP46215.2023.10179435}

\bibitem{poisonedrag2024}
Z.\ Zou, R.\ Geng, B.\ Wang, H.\ Chen, and N.\ Zhang,
``PoisonedRAG: Knowledge Corruption Attacks to Retrieval-Augmented Generation
of Large Language Models,''
\textit{arXiv preprint arXiv:2402.07867}, 2024.
\newblock DOI: \texttt{10.48550/arXiv.2402.07867}

\bibitem{plonk2019}
A.\ Gabizon, Z.\,J.\ Williamson, and O.\ Ciobotaru,
``PLONK: Permutations over Lagrange-Bases for Oecumenical Noninteractive
Arguments of Knowledge,''
\textit{IACR Cryptology ePrint Archive}, Report 2019/953, 2019.
\newblock URL: \texttt{https://eprint.iacr.org/2019/953}

\bibitem{peng2025survey}
Z.\ Peng, C.\ Zhao, T.\ Wang, G.\ Liao, Z.\ Lin, Y.\ Liu, B.\ Cao,
L.\ Shi, Q.\ Yang, and S.\ Zhang,
``A Survey of Zero-Knowledge Proof Based Verifiable Machine Learning,''
\textit{Artificial Intelligence Review}, Springer, 2026.
\newblock DOI: \texttt{10.1007/s10462-026-11557-y}

\bibitem{defitrace2024}
DeFiTrace: Event-Enriched Detection of Price Oracle Manipulation Across DeFi
Transactions,
\textit{ACM Trans.\ Privacy and Security}, 2024.
\newblock DOI: \texttt{10.1145/3817054}

\bibitem{bensasson2018stark}
E.\ Ben-Sasson, I.\ Bentov, Y.\ Horesh, and M.\ Riabzev,
``Scalable, Transparent, and Post-Quantum Secure Computational Integrity,''
\textit{IACR Cryptology ePrint Archive}, Report 2018/046, 2018.
\newblock URL: \texttt{https://eprint.iacr.org/2018/046}

\bibitem{karim2025aibc}
M.\,M.\ Karim, D.\,H.\ Van, S.\ Khan, Q.\ Qu, and Y.\ Kholodov,
``AI Agents Meet Blockchain: A Survey on Secure and Scalable Collaboration
for Multi-Agents,''
\textit{Future Internet}, vol.\ 17, no.\ 2, p.\ 57, MDPI, 2025.
\newblock DOI: \texttt{10.3390/fi17020057}

\bibitem{sokaibc2025}
SoK: Security and Privacy of AI Agents for Blockchain,
\textit{arXiv preprint arXiv:2509.07131}, accepted BCCA 2025.
\newblock DOI: \texttt{10.48550/arXiv.2509.07131}

\bibitem{defort2024}
Y.\ Xie, S.\ Zhang, J.\ Liu, X.\ Zhao, and H.\ Wang,
``DeFort: Automatic Detection and Analysis of Price Manipulation Attacks in
DeFi Applications,''
in \textit{Proc.\ ISSTA 2024}, ACM, 2024.
\newblock DOI: \texttt{10.1145/3650212.3652137}

\end{thebibliography}

\end{document}
